\documentclass[a4paper,11pt]{ltjsarticle}
% 数式
\usepackage{amsmath,amsfonts}
\usepackage{bm}
% 画像
\usepackage{graphicx}
\usepackage{circuitikz}
\usepackage{amsmath,amssymb}
\usepackage{siunitx}
\usepackage{float}
\usepackage{tikz}
\usepackage{askmaps}
\usepackage{multirow}
\usepackage{bigstrut}
\usepackage{slashbox}
\usepackage{rotating}
\usepackage{listings}
% 数式
\usepackage{physics}
\usepackage{mathtools}
% 画像
\usepackage{subcaption}
% 表
\usepackage{makecell}
% その他
\usepackage{url}
\usepackage{ascmac}
\usepackage{cases}
\usepackage{here}
\usepackage{upgreek}
% 日本語対応
\usepackage{luatexja}
\usepackage{luatexja-fontspec}

\AtBeginDocument{\RenewCommandCopy\qty\SI}


\definecolor{commentgreen}{RGB}{0,200,0}
\definecolor{eminence}{RGB}{120,80,250}
\definecolor{weborange}{RGB}{255,165,0}
\definecolor{frenchplum}{RGB}{10,150,200}
\definecolor{commentgreen}{RGB}{0,200,0}
\definecolor{eminence}{RGB}{120,80,250}
\definecolor{weborange}{RGB}{255,165,0}
\definecolor{frenchplum}{RGB}{10,150,200}

\lstset{
        language = {C},
        basicstyle = \ttfamily\small,
        keywordstyle=\color{eminence}\ttfamily\bfseries,
        commentstyle=\color{commentgreen}\textit,
    identifierstyle=\color{black}\ttfamily,
        xleftmargin=.35in,
        frame=lines,
    showstringspaces=false,
        numbers=left,
        stepnumber = 1,
        breaklines=true,
        numberstyle = \ttfamily\normalsize,
    tabsize=4,  
        emph={int, int8_t, int16_t, int32_t, int64_t, uint8_t, uint16_t, uint32_t, uint64_t, char, double, float, unsigned, void, bool},
        emphstyle={\color{blue}}, 
        morekeywords={>, <, ., ;, +, -, *, /, !, =, ~},
        breakindent = 10pt, 
        framexleftmargin=10mm, 
        columns=fixed,
        basewidth=0.5em,
        }

% 特定のスタイル設定
\lstdefinestyle{customtxt}{
  basicstyle=\ttfamily\footnotesize,
  backgroundcolor=\color{lightgray},
  frame=single,
  breaklines=true,
  columns=fullflexible,
  showspaces=false,
  showstringspaces=false,
  showtabs=false,
  tabsize=4,
}

\newcommand{\fig}[4]{
    \begin{figure}[htbp]
      \centering
      \includegraphics{./image/#1}
      \caption{#2}
      \label{fig:#3}
    \end{figure}
  }
\begin{document}


\section{要旨}
% 結果等を入れて三行程度でミニレポートを書く
物質を曲げようと応力を加えると，ひずみが発生する．応力$\sigma$とひずみ$\varepsilon$の関係は$\sigma = E \varepsilon$で表される．
ここで，$E$はヤング率と呼ばれる物質固有の定数である．
$E$はフックの法則に現れる$k$に相当する．
本実験では，様々な金属を用いてヤング率の計測を行い，実測値との比較を行った．
銅のヤング率が約$12.2 \times 10^{10}$,
黄銅のヤング率が約$9.41 \times 10^{10}$,
鋼のヤング率が約$20.6 \times 10^{10}$であった．
\section{目的}
% 何を測定して、どのような意味があるのかを書く
金属棒におもりを加えていき，ひずみを測定することにより，ヤング率を求める．
また，これらを3種類の金属棒で行い，それぞれのヤング率を求める．
\section{実験手順}
% 実験指導書のページを参照で良い。しかし違うものをした場合は再現できるほど細かく書く
本実験は実験指導書pp.37-43を参考にして行う．
\section{実験結果}
% 表/グラフを用いて実験で得たデータを示し、目的とする物理量を求める。最低限有効数字を意識して、さらには不確かさを表示すること
\subsection{測定データ}
実験で得られた測定データを示す．
金属棒は銅, 黄銅, 鋼の3種類についてそれぞれ，A, B, Cとする．
金属棒A,B,Cそれぞれの測定時のエッジ間の距離$l$を表\ref{tab:metal_length}に示す．
金属棒A,B,Cのそれぞれの厚さ，長さ，幅を表\ref{tab:metal_size}に示す．
金属棒Aのおもりを加えていった時の表を表\ref{tab:youngs_modulus_A}に示す．
金属棒Bのおもりを加えていった時の表を表\ref{tab:youngs_modulus_B}に示す．
金属棒Cのおもりを加えていった時の表を表\ref{tab:youngs_modulus_C}に示す．
\begin{table}[htbp]
  \centering
  \caption{金属棒のエッジ間の距離}
  \label{tab:metal_length}
  \begin{tabular}{|c|c|}
    \hline
    金属 & $x_m$[cm] \\ \hline
    A  & 240.5     \\ \hline
    B  & 237       \\ \hline
    C  & 234       \\ \hline
  \end{tabular}
\end{table}
\begin{table}[htbp]
  \centering
  \caption{金属棒のサイズ}
  \label{tab:metal_size}
  \begin{tabular}{|c|c|c|c|}
    \hline
    金属棒 & 幅b$[mm]$ & 厚さa$[mm]$ & 断面積$[mm^2]$ \\ \hline
    A   & 15.870   & 4.9700    & 78.87       \\ \hline
    B   & 15.959   & 4.9470    & 78.95       \\ \hline
    C   & 15.872   & 4.8965    & 77.72       \\ \hline
  \end{tabular}
\end{table}
\begin{table}[htbp]
  \centering
  \caption{金属棒Aの測定結果}
  \label{tab:youngs_modulus_A}
  \begin{tabular}{|c|c|c|c|}
    \hline
    $m$ & $y_m$[mm] & $y'_m$[mm] & $\overline{y}$[mm] \\
    \hline
    0   & 0.0       & -0.5       & -0.3               \\
    \hline
    1   & 20.5      & 20.3       & 20.4               \\
    \hline
    2   & 41.2      & 41.5       & 41.4               \\
    \hline
    3   & 62.3      & 61.7       & 62.0               \\
    \hline
    4   & 83.5      & 83.5       & 83.5               \\
    \hline
    5   & 104.5     & 104.5      & 104.5              \\
    \hline
  \end{tabular}
\end{table}
\begin{table}[htbp]
  \centering
  \caption{金属棒Bの測定結果}
  \label{tab:youngs_modulus_B}
  \begin{tabular}{|c|c|c|c|}
    \hline
    $m$ & $y_m$[mm] & $y'_m$[mm] & $\overline{y}$[mm] \\ \hline
    0   & 30.1      & 34.0       & 32.1               \\ \hline
    1   & 60.3      & 60.5       & 60.4               \\ \hline
    2   & 82.3      & 87.5       & 84.9               \\ \hline
    3   & 114.5     & 114.5      & 114.5              \\ \hline
    4   & 141.0     & 141.0      & 141.0              \\ \hline
    5   & 167.5     & 167.5      & 167.5              \\ \hline
  \end{tabular}
\end{table}
\begin{table}[htbp]
  \centering
  \caption{金属棒Cの測定結果}
  \label{tab:youngs_modulus_C}
  \begin{tabular}{|c|c|c|c|}
    \hline
    $m$ & $y_m$[mm] & $y'_m$[mm] & $\overline{y}$[mm] \\ \hline
    0   & 30.0      & 30.0       & 30.0               \\ \hline
    1   & 41.5      & 42.0       & 41.8               \\ \hline
    2   & 55.3      & 55.0       & 55.2               \\ \hline
    3   & 67.5      & 68.5       & 68.0               \\ \hline
    4   & 81.5      & 81.0       & 81.3               \\ \hline
    5   & 94.0      & 94.0       & 94.0               \\ \hline
  \end{tabular}
\end{table}
\newpage
\subsection{ヤング率の計算}
ヤング率は次の式で求めることができる．
\begin{gather}
  e = \frac{z \cdot \overline{y}}{2x} \\
  E = \frac{l^3mg}{4a^3be}
\end{gather}
測定結果より求めた


金属棒A,B,Cのヤング率を表\ref{tab:youngs_modulus}に示す．
\begin{table}[htbp]
  \centering
  \caption{ヤング率の計算結果}
  \label{tab:youngs_modulus}
  \begin{tabular}{|c|c|c|c|}
    \hline
    金属棒 & ヤング率$E$[$\mathrm{N/m^2}$] & ヤング率$E$[$\mathrm{Pa}$] & ヤング率$E$[$\mathrm{GPa}$] \\ \hline
    A   & 1.22$\times 10^{11}$      & 1.22$\times 10^{11}$   & 122                     \\ \hline
    B   & 9.41$\times 10^{10}$      & 9.41$\times 10^{10}$   & 94.1                    \\ \hline
    C   & 2.06$\times 10^{11}$      & 2.06$\times 10^{11}$   & 206                     \\ \hline
  \end{tabular}
\end{table}
\section{考察}
% 実験結果から何がいえるかという問に答えるところ。実験結果が信頼できるのかを論理的に書く。実験指導書の課題もここで答える。新しい実験方法を披露できるとなおよい。最小二乗法は考察
理論値の値を表\ref{tab:youngs_modulus_true}に示す．実験値と理論値の誤差は少なかったが，
誤差は生じていた．この誤差の生じる原因として，エッジ間の距離の測定誤差が考えられる．
鏡の向きを合わせた後に測定を行ったため，目測で大体あっていることを確認しながらの計測，またメジャーを平行にすることの難しさから
この誤差が生じていたと思われる．
誤差が少なかった要因としては，3乗される数値の測定誤差の少なさが挙げられる．$b$,$a$の測定箇所による数値の誤差は0.1mm程度であり，どの測定箇所でもほぼ同じ値となった．これにより
実測値と理論値の誤差が少なくなったと考えられる．
\begin{table}[htbp]
  \centering
  \caption{実際のヤング率}
  \label{tab:youngs_modulus_true}
  \begin{tabular}{|c|c|}
    \hline
    金属棒 & ヤング率$E$[$\mathrm{N/m^2}$]   \\ \hline
    銅   & $12.98\times 10^{10}$       \\ \hline
    黄銅  & $10.06\times 10^{10}$       \\ \hline
    鋼   & $20.1 ~ 20.6\times 10^{10}$ \\ \hline
  \end{tabular}
\end{table}




\end{document}